% SPDX-FileCopyrightText: Copyright (c) 2026 Yegor Bugayenko
% SPDX-License-Identifier: MIT

\documentclass{article}
\usepackage{hyperref}
\usepackage{verbatim}

\title{Quality Control in Software Projects}
\author{Anonymous}
\date{2026}

\begin{document}

\maketitle

\begin{abstract}
Quality is non-negotiable.
Yet most teams delegate the responsibility for it to the same
  programmers who are paid to ship features quickly, which produces
  a conflict of interest that the project rarely wins.
This paper argues that quality must be enforced by infrastructure,
  not by goodwill.
\end{abstract}

\section{Introduction}

Speed and quality pull in opposite directions.
A programmer paid to close tickets has every reason to cut corners,
  and a programmer paid to maintain quality has every reason to slow
  down.
When the same person is asked to do both, one of the two goals
  always loses.

The mistake is asking.

A mature project does not ask programmers to police themselves.
It builds the rules into the pipeline, the repository, and the
  review process, so that low-quality work cannot enter the codebase
  even when someone tries to push it through.

\section{The Programmer's Incentive}

Programmers ship.
That is the work.

A team that rewards velocity should not be surprised when its
  members optimise for velocity, including by skipping tests,
  ignoring lint warnings, and merging without review.
Asking them to defend quality on top of shipping is asking them to
  fight their own paycheque.

\section{The Infrastructure's Job}

Quality belongs to the infrastructure.
The build system, the version control rules, and the review pipeline
  must reject bad work without human judgement.

The minimum bar for a serious project includes:

\begin{itemize}
  \item A read-only master branch that accepts only merged pull
    requests.
  \item A pre-merge build that runs every test, every linter, and
    every static analyzer.
  \item Coverage thresholds that fail the build when they slip.
  \item At least two independent reviewers on every change.
  \item A rule that no one merges their own code, ever.
\end{itemize}

None of these rules trust the author of the change.
Each of them assumes that the author will, at some point, try to
  slip something through, and each closes one route by which they
  could succeed.

\section{A Concrete Example}

Suppose the build invokes the linter through a wrapper script.
The wrapper must fail the build on any warning, not just on errors,
  because a warning today is an error tomorrow.

\begin{verbatim}
$ ./scripts/lint.sh --fail-on-warning --report=junit --output=build/lint-report.xml
\end{verbatim}

A line like that one is ugly to read but trivial to enforce, and
  enforcement is the point.
The build either passes or it does not.
No human reads the output and decides whether the warning matters.

\section{The Caveat}

This model only works in projects mature enough to honour it.
A team that disables the linter under deadline pressure has no
  infrastructure, only the appearance of one.
Worse, a team that allows ``temporary'' bypasses creates a folklore
  in which every bypass is temporary and none are ever removed.

The discipline is binary.
Either the rules hold for everyone, every time, or they do not
  exist.

\section{Conclusion}

Programmers should be fast.
Infrastructure should be strict.
When those two forces are aligned against each other, the project
  ships quickly without losing the quality it cannot afford to lose.

\end{document}
